\documentclass[10pt]{letter}
\usepackage[T1]{fontenc}
\usepackage[utf8]{inputenc}
\usepackage[margin=0.8in]{geometry}
\usepackage[hidelinks]{hyperref}

\signature{Taiqian Liu\\on behalf of Lijun Liu and Chunhui Zhu}
\address{Department of Automation\\
School of Aerospace Engineering\\
Xiamen University\\
Xiamen 361005, China\\
\href{mailto:liulijun@xmu.edu.cn}{liulijun@xmu.edu.cn}}

\begin{document}
\begin{letter}{Editor\\
\textit{Computer Methods in Applied Mechanics and Engineering}}

\opening{Dear Editor,}

We submit the manuscript ``Boundary-Capacitance Neural Operators for High-Contrast Multiscale Homogenization'' for consideration in \textit{Computer Methods in Applied Mechanics and Engineering}.

Surrogate modeling of representative volume elements (RVEs) is essential to two-scale $\mathrm{FE}^2$ homogenization, yet standard neural operators suffer $\Theta(\kappa)$ error amplification at high phase contrast, while equipotential clamping discards the $O(1)$ interior polarization flux. Perfect-conductor asymptotics work at $\kappa\ge 10^4$; the practical gap is the intermediate band $\kappa\in[10,10^3]$. We characterize the discrete sensitivity $\|g\|_1=\Theta(\kappa)$, prove one-cell strip localization of $\sigma=h L_0(u-\ell)$, and relocate standard $\mathrm{FE}^2$ virtual-work averaging from the outer RVE boundary to the inclusion interface, yielding contrast-invariant readout on predicted fields. The Boundary-Capacitance Neural Operator (BC-NO) couples $O(N^d\log N)$ constant-coefficient Green layers with normalized residuals, capacitance pooling, and percolation-aware adaptive readouts.

Under a shared dipole readout across fourteen benchmarks, architectural OOD gains versus the strongest baseline per experiment are $1.0\times$--$4.1\times$ (readout often explains the rest). Thermal and elastic $\mathrm{FE}^2$ plates reach $2.38\%$ flux and $0.61\%$ displacement error. Same-GPU Jacobi-PCG speedups are $2.7\times$--$8.0\times$; same-CPU sparse LU is slower at $N=128$ ($0.86\times$). GPU-vs-CPU $\mathrm{FE}^2$ wall-clock ratios are reported as workflow comparisons. Limitations include percolating fallback, cross-resolution scaling, $32^3$ synthetic accuracy, and the absence of matched-tolerance AMG/FFT-CG timing.

The code and data generators to reproduce all results will be hosted at \url{https://github.com/LTQ12/BC-NO}. This manuscript is original, has not been published previously, and is not under consideration elsewhere. All authors have approved the manuscript and its submission.

Thank you for considering our work.

\closing{Sincerely,}

\end{letter}
\end{document}
